%% ============================================================
%%  Capítulo 11 – Trazabilidad de Requisitos
%% ============================================================
\chapter{Trazabilidad de Requisitos}
\label{chap:trazabilidad}

\section{Metodología}
\label{sec:traz_metodologia}

La matriz de trazabilidad se construyó siguiendo una cadena directa: \textbf{requisito del E1 → decisión de diseño del E3 → componente o bloque que la implementa → evidencia documentada en este informe}. Para cada requisito se indica la clasificación (HW puro, SW puro, o Mixto HW+SW) y el método de verificación que podrá aplicarse en el prototipo físico del Entregable 4.

\section{Matriz de Trazabilidad Requisito--Diseño}
\label{sec:traz_matriz}

\begin{longtable}{L{1.5cm} L{3cm} C{1.3cm} L{3.2cm} L{3cm} C{1.5cm}}
  \caption{Matriz de trazabilidad requisito--diseño de ALMA}
  \label{tab:trazabilidad} \\
  \toprule
  \rowcolor{udea-green}
  \textcolor{white}{\textbf{ID}} &
  \textcolor{white}{\textbf{Descripción}} &
  \textcolor{white}{\textbf{Clasif.}} &
  \textcolor{white}{\textbf{Comp. / Bloque}} &
  \textcolor{white}{\textbf{Decisión clave}} &
  \textcolor{white}{\textbf{Sección}} \\
  \midrule
  \endfirsthead
  \toprule
  \rowcolor{udea-green}
  \textcolor{white}{\textbf{ID}} &
  \textcolor{white}{\textbf{Descripción}} &
  \textcolor{white}{\textbf{Clasif.}} &
  \textcolor{white}{\textbf{Comp. / Bloque}} &
  \textcolor{white}{\textbf{Decisión clave}} &
  \textcolor{white}{\textbf{Sección}} \\
  \midrule
  \endhead
  \bottomrule
  \endfoot
  \reqref{REQ-F-01} & Reproducción audio HD 24-bit/96\,kHz & Mixto &
    Amp. TPA3118, Códec PCM5122, bus I2S del A311D2 &
    Bus I2S 24-bit + Amp. Clase D; riel VCC5 de pureza espectral &
    \S\ref{sec:arq_audio}, \S\ref{sec:ctrl_bloques} \\
  \rowcolor{row-alt}
  \reqref{REQ-F-02} & Captura e interacción por voz & Mixto &
    Micros MEMS SPH0645 (×2), PDM nativo A311D2, OVOS &
    Micrófonos MEMS PDM dual; procesamiento local en A311D2 &
    \S\ref{sec:arq_audio}, \S\ref{sec:ctrl_bloques} \\
  \reqref{REQ-F-03} & Persistencia de configuración ante corte & Mixto &
    eMMC Samsung 32\,GB, Flash QSPI W25Q256, STM32 &
    eMMC soldada con TPS2115; ext4 con journaling &
    \S\ref{sec:ctrl_bloques}, \S\ref{sec:hwsw_arquitectura} \\
  \rowcolor{row-alt}
  \reqref{REQ-F-04} & Iluminación RGB dinámica & HW+SW &
    18 LEDs WS2812B, STM32G030 (PMU) &
    STM32 controla cadena WS2812 vía GPIO; recibe comandos I2C del SoC &
    \S\ref{sec:ctrl_bloques}, \S\ref{sec:hwsw_hw_condiciones} \\
  \reqref{REQ-NF-01} & Silenciado físico de micrófonos & HW &
    Switch DPST JS202011, LED rojo D30 &
    Corte galvánico hard-wired en riel alimentación micrófonos &
    \S\ref{sec:ctrl_bloques}, \S\ref{sec:hwsw_flujos} \\
  \rowcolor{row-alt}
  \reqref{REQ-NF-02} & Control de dispositivos domóticos vía IP & SW &
    Módulo WiFi/BT AP6256, A311D2 &
    ALMA como cliente IP; TLS\,1.3 gestionado por TrustZone A311D2 &
    \S\ref{sec:arq_conectividad}, \S\ref{sec:hwsw_arquitectura} \\
  \reqref{REQ-G-01} & GUI táctil de baja latencia & Mixto &
    Panel MIPI DSI (4 lanes), controlador táctil GT911 &
    Interfaz MIPI DSI 4L del A311D2; driver DRM/KMS en kernel Linux &
    \S\ref{sec:ctrl_bloques}, \S\ref{sec:hwsw_arquitectura} \\
  \rowcolor{row-alt}
  \reqref{REQ-A-01} & Gestión térmica -- T\textsubscript{j} $<$ 105\,°C & HW &
    PCB-CTRL (copper pours bajo SoC), PCB-PWR (pads térmicos reguladores) &
    Partición en 2 PCBs; thermal vias bajo BGA; disipadores en reguladores &
    \S\ref{chap:particion}, \S\ref{chap:fabricante} \\
  \reqref{REQ-E-01} & Seguridad eléctrica IEC\,62368-1 & HW &
    Fusible PPTC MF-MSMF300, TVS SMAJ15CA, Power mux TPS2121 &
    Fusible PPTC + TVS en entrada 12\,V; clearance $\geq$0.5\,mm a 12\,V &
    \S\ref{sec:pwr_bloques}, \S\ref{chap:fabricante} \\
  \rowcolor{row-alt}
  \reqref{REQ-E-02} & Ciberseguridad ETSI EN\,303\,645 & Mixto &
    TrustZone integrado A311D2, Flash QSPI (bootloader firmado) &
    TrustZone para arranque seguro y gestión de claves; U-Boot con firma &
    \S\ref{sec:arq_seguridad}, \S\ref{sec:hwsw_flujos} \\
\end{longtable}

\section{Análisis de Cobertura}
\label{sec:traz_cobertura}

La Tabla~\ref{tab:trazabilidad} demuestra que \textbf{los diez requisitos del E1 tienen trazabilidad completa} hacia componentes concretos, decisiones de diseño documentadas y secciones verificables en este informe. No existe ningún requisito sin componente asignado ni decisión de diseño indefinida al cierre del E3.

Los requisitos de naturaleza mixta HW+SW (REQ-F-01, REQ-F-02, REQ-F-03, REQ-G-01, REQ-E-02) tienen la parte de hardware completamente definida en este entregable. La validación de la parte de software se realizará en la fase de integración, fuera del alcance del proyecto académico, pero el hardware dimensionado provee la base técnica necesaria para su implementación.

\section{Requisitos con Implementación Pendiente o Parcial}
\label{sec:traz_pendientes}

No existen requisitos con implementación de hardware pendiente al cierre del E3. Todos los bloques y componentes han sido seleccionados con número de parte concreto y sus restricciones de diseño han sido trasladadas a las reglas DRC y SI de Altium Designer.

El único aspecto que permanece fuera del alcance de este entregable, por diseño, es la verificación física de los requisitos sobre el prototipo (pruebas de THD+N del amplificador, latencia de wake-word, temperatura de unión del SoC bajo carga), que corresponde a la fase de prototipado posterior al Entregable 4.
